\documentclass{article}
\usepackage[backend=biber, style=authoryear, sorting=nyt, style=numeric]{biblatex}
\usepackage{url}
\usepackage[dvipsnames]{xcolor}
\usepackage{enumitem}

\usepackage{background}
\usepackage{tikz}
\usepackage{geometry}
\usetikzlibrary{shapes, arrows, positioning}
\usepackage[T1]{fontenc}
\usepackage[utf8]{inputenc}
\usepackage[british]{babel}
\usepackage{titletoc}
\usepackage{titlesec}
\usepackage{lipsum}
\usepackage{tocloft}
\usepackage{hyperref}
\usepackage{algorithm}
\usepackage{algpseudocode}
\usepackage{graphicx}
\usepackage[section]{placeins}
\usepackage{caption}
\usepackage{textcomp}
\usepackage{subcaption}
\usepackage[edges]{forest}
%\usepackage[superscript,biblabel]{cite}
\usepackage{fontawesome5}
\usepackage{algorithm}
\usepackage{algpseudocode}
\usepackage{amsmath}
\usepackage[simplified]{pgf-umlcd}
\usepackage{enumitem}
\usepackage{adjustbox}
\usepackage{listings}
\usetikzlibrary{shapes.geometric, arrows}

\tikzset{
	startstop/.style={
		rectangle,
		rounded corners,
		minimum width=3cm,
		minimum height=1cm,
		text centered,
		draw=black,
		fill=red!30
	},
	io/.style={
		trapezium,
		trapezium stretches=true, % A later addition
		trapezium left angle=70,
		trapezium right angle=110,
		minimum width=3cm,
		minimum height=1cm,
		text centered,
		draw=black,
		fill=blue!30
	},
	process/.style={
		rectangle,
		minimum width=3cm,
		minimum height=1cm,
		text centered,
		text width=3cm,
		draw=black,
		fill=orange!30
	},
	decision/.style={
		diamond,
		minimum width=3cm,
		minimum height=1cm,
		text centered,
		draw=black,
		fill=green!30
	},
	arrow/.style={
		thick,
		->,
		>=stealth
	}
}

\renewcommand{\cfttoctitlefont}{\hfill\large\bfseries\fontsize{20}{24}\selectfont}
\renewcommand{\cftaftertoctitle}{\hfill\mbox{}}

\definecolor{mydarkyellow}{HTML}{87CEEB}

\backgroundsetup{
	scale=1,
	color=black,
	opacity=0.5,
	angle=0,
	contents={
		\ifnum\value{page}=1
		\begin{tikzpicture}[remember picture,overlay]
			\path [fill=mydarkyellow] (current page.south west) rectangle (current page.north east);
		\end{tikzpicture}
		\fi
	}
}

\title{\textbf{\Huge Polytech Nice}\\[1cm]
	\textbf{\LARGE Project report}\\[2cm]
	\hrule height 1pt
	\vspace{0.5cm}
	\textbf{\Large numerical solution of the optimal orbit transfer problem in the 
		plane}\\[0.5cm]
	\hrule height 1pt
	\vspace{3cm}
	\small{\today}}

\author{
	\begin{tabular}{c}
		Gerbaud Florent \\
	\end{tabular}
}

\date{}

\begin{document}
	\maketitle
	\vspace{3cm}
	\begin{center}
		\includegraphics[width=5cm]{Polytech.png}
	\end{center}
	\newpage % ajout d'un saut de page
	% Table des matières
	\renewcommand{\contentsname}{
		\hfill
		\begin{tikzpicture}
			\node[draw, fill=white, inner sep=20pt,line width=1.5pt] {\fontsize{30}{36}\selectfont\bfseries Table des matières};
		\end{tikzpicture}
		\hfill
	}
	\tableofcontents
	\listofalgorithms
	\listoffigures
	\newpage
	
	\section{Theory}
	
	\subsection{Nomenclature}
		\begin{itemize}
			\item r := satellite radial distance attracting body
			\item u := radial velocity
			\item v := tangential speed
			\item m := satellite mass
			\item $\frac{\partial m}{\partial t} = \dot{m} = -\frac{T}{g_0 \cdot Isp} := $  dm motor mass flow
			\item $\phi := $ angle of thrust direction with r
			\item $\mu := $  gravitational constant of the attracting body
			\item $\theta := $ the angle formed by r with r(0)
		\end{itemize}
		
		\subsection{Conditions initiales et finales:}
		\[ \boxed{\begin{aligned}
				& r(0)=r_0, \quad u(0)=0, \quad v(0)=\sqrt{\frac{\mu}{r_0}} \\
				& r\left(t_f\right)=r_f, \quad u\left(t_f\right)=0, \quad v\left(t_f\right)=\sqrt{\frac{\mu}{r_f}}
		\end{aligned}} \]
	 
		
		\subsection{Hamiltonien et dynamique de l'état adioint:}
		\begin{align*}
			& \dot{x}=\begin{bmatrix}
				\dot{r} \\
				\dot{u} \\
				\dot{v}
			\end{bmatrix}=f(x, t, \phi)=\begin{bmatrix}
				u \\ \\
				\frac{v^2}{r}-\frac{\mu}{r^2}+\frac{T}{m_0+\dot{m} \times t} \sin (\phi) \\ \\
				-\frac{u v}{r}+\frac{T}{m_0+\dot{m} \times t} \cos (\phi)
			\end{bmatrix} \\ \\
			& \lambda^T=\begin{bmatrix}
				\lambda_r & \lambda_u & \lambda_v
			\end{bmatrix} \\ \\
			& H=1+\lambda^T f=1+\lambda_r u+\lambda_u\left(\frac{v^2}{r}-\frac{\mu}{r^2}+\frac{T}{m_0+\dot{m} \times t} \sin (\phi)\right)+\lambda_v\left(-\frac{u v}{r}+\frac{T}{m_0+\dot{m} \times t} \cos (\phi)\right) \\ \\
			& \dot{\lambda}=\begin{bmatrix}
				\dot{\lambda}_r \\
				\dot{\lambda}_u \\
				\dot{\lambda}_v
			\end{bmatrix}=-\frac{\partial f^T}{\partial x} \lambda=\begin{bmatrix}
				-\lambda_u\left(-\frac{v^2}{r^2}+\frac{2 \mu}{r^3}\right)-\lambda_v\left(\frac{u v}{r^2}\right) \\ \\
				-\lambda_r+\lambda_v \frac{v}{r} \\ \\
				-\lambda_u \frac{2 v}{r}+\lambda_v \frac{u}{r}
			\end{bmatrix}
		\end{align*}
		
		\subsection{Equation de la commande}
		\begin{align*}
			& H=1+\lambda^T f=1+\lambda_r u+\lambda_u\left(\frac{v^2}{r}-\frac{\mu}{r^2}+\frac{T}{m_0+\dot{m} \times t} \sin (\phi)\right)+\lambda_v\left(-\frac{u v}{r}+\frac{T}{m_0+\dot{m} \times t} \cos (\phi)\right) \\ \\
			& \frac{\partial H}{\partial \phi}=\lambda_u\left(\frac{T}{m_0+\dot{m} \times t} \cos (\phi)\right)-\lambda_v\left(\frac{T}{m_0+\dot{m} \times t} \sin (\phi)\right)=0 \\ \\
			& \frac{\partial H}{\partial \phi}=0 \Leftrightarrow \lambda_u \cos (\phi)-\lambda_v \sin (\phi)=0 \Rightarrow \tan (\phi)=\frac{\lambda_u}{\lambda_v} \\
		\end{align*}
		
		\subsection{Condition d'optimalité sur $t_f$}
		\[
		\\
		\left[H=1+\lambda^T f=1+\lambda_r u+\lambda_u\left(\frac{v^2}{r}-\frac{\mu}{r^2}+\frac{T}{m_0+\dot{m} \times t} \sin (\phi)\right)+\lambda_v\left(-\frac{u v}{r}+\frac{T}{m_0+\dot{m} \times t} \cos (\phi)\right)\right]_{t_j}=0
		\\ \] \newline
		
		Cette condition peut être remplacée par la condition équivalente à $t_0=0$:
		\[
		\left[\lambda^T \times \lambda\right]_{t_0}=1
		\]
		
		\subsection{Formulation du problème aux deux bouts}
		Inconnues à $t=t_0=0$: $\lambda(t_0)^T = [\lambda_r \ \lambda_u \ \lambda_v]_{t_0}$ \\
		Inconnue à $t=t_f$: $t_f$
		Système d'équations à résoudre:
		\begin{align*}
			g(\lambda_{r0}, \lambda_{u0}, \lambda_{v0}, t_f) &= \left[\begin{array}{c}
				r(t_f) - r_f \\ \\
				u(t_f) \\ \\
				v(t_f) - \sqrt{\frac{\mu}{r_f}} \\ \\
				\left[1 + \lambda^T f\right]_{t_f}
			\end{array}\right] = 0
		\end{align*}
		
	\section{Implementation}
	
	\begin{algorithm}[H]
		\caption{Problem data declaration}\label{alg:satellite_maneuvering}
		\begin{algorithmic}
			%\State \textbf{Input:} \\
			%\begin{itemize}
			%\State AU = 149597870690; // 1 astronomical unit in m
			%\State $r0 = AU$; // Initial orbit radius in m
			%\State $rf = 1.5 \times AU$; // Final target orbit radius in m
			%\State $m0 = 1000$; // Satellite mass at t0 in kg
			%\State $T = 0.1:0.1:0.6$; // Engine thrust in N (creates a vector from 0.1 to 0.6 with a step of 0.01)
			%\State $g0 = 9.80665$; // Earth gravity in ms^-2
			%\State $Isp = 3000$; // Specific impulse of the engine in s
			%\State $\mu_{\text{body}} = 1.32712440018 \times 10^{20}$; // Gravitational constant of the sun in m^3s^-2
			%\end{itemize}
			\State \textbf{Data:} \\
			\begin{itemize}[]
				\item \textbf{\underline{Initial Conditions:}}
				\begin{itemize}
					\item $AU = 149597870690$
					\item $r_0 = AU$
					\item $u_0 = 0$
					\item $v_0 = \sqrt{\frac{\mu_{\text{body}}}{r_0}}$
					\item$ m_0=1000$
				\end{itemize}
				\item \textbf{\underline{Final Conditions:}}
				\begin{itemize}
					\item $r_f = 1.5 \times AU$
					\item $u_f = 0$
					\item $v_f = \sqrt{\frac{\mu_{\text{body}}}{r_f}}$
				\end{itemize}
				\item \textbf{\underline{Power}}:
				\begin{itemize}
					\item Set $T$ to an array containing values from $0.1$ to $0.6$ with a step size of $0.1$.
				\end{itemize}
				\item \textbf{\underline{Earth gravity}}
				\begin{itemize}
					\item $g_0 = 9.80665$
				\end{itemize}
				\item \textbf{\underline{Specific motor impulse}}
				\begin{itemize}
					\item $Isp = 3000$
				\end{itemize}
				\item \textbf{\underline{sun's gravitational constant}}
				\begin{itemize}
					\item $\mu_{body} = 1.32712440018e+20$
				\end{itemize}
			\end{itemize}
		\end{algorithmic}
	\end{algorithm}
	
	
	\subsection{Problem Data declaration}
	
	In this code, we have to normalize the data due to the problem of orders of magnitude. This allows us to make the resolution easier and more accurate for the solver.
	
	\begin{algorithm}[H]
		\caption{Problem data declaration}\label{alg:satellite_maneuvering}
		\begin{algorithmic}
			
			%\end{itemize}
			\State \textbf{Normalized unit:} \\
			\begin{itemize}[]
				\item \textbf{\underline{Unit:}}
				\begin{itemize}
					\item $DU = r_0$; % Distance unitaire
					\item $VU = v_0$; % Vitesse unitaire
					\item $MU = m_0$; % Masse unitaire
					\item $TU = \frac{DU}{VU}$; % Temps unitaire
					\item $FU = \frac{MU \times DU}{TU^2}$; %
				\end{itemize}
			\end{itemize}
		\end{algorithmic}
	\end{algorithm}
	

	\newpage
	\subsection{System Dynamics Resolution function}	
	\begin{algorithm}[H]
		\caption{Dynamical model of a spacecraft}
		\label{alg:dynamical_model}
		\begin{algorithmic}
			\State \textbf{\underline{Definition of the struct:}}
			\begin{algorithmic}
				\State $\text{param}.DU = DU;$
				\State $\text{param}.VU = VU;$
				\State $\text{param}.MU = MU;$
				\State $\text{param}.TU = TU;$
				\State $\text{param}.FU = FU;$
				\State $\text{param}.\mu_{\text{body}} = \frac{\mu_{\text{body}}}{\text{param.DU}^3 \times \text{param.TU}^2};$
				\State $\text{param}.T = \frac{T(1)}{\text{param.FU}};$
				\State $\text{param}.g0_{\text{Isp}} = \frac{g0 \times Isp}{\text{param.VU}};$
				\State $\text{param}.m0 = \frac{m0}{\text{param.MU}};$
				\State $\text{param}.t0 = 0;$
				\State $\text{param}.x_{t0} = \begin{bmatrix} \frac{r0}{\text{param.DU}} \\ \frac{u0}{\text{param.VU}} \\ \frac{v0}{\text{param.VU}} \end{bmatrix};$
				\State $\text{param}.x_{tf} = \begin{bmatrix} \frac{rf}{\text{param.DU}} \\ \frac{uf}{\text{param.VU}} \\ \frac{vf}{\text{param.VU}} \end{bmatrix};$
			\end{algorithmic}
			
		\end{algorithmic}
	\end{algorithm}
	
	\begin{algorithm}[H]
		\caption{Function \texttt{dynpol}}
		\label{alg:dynpol}
		\begin{algorithmic}
			\State \textbf{Input:}
			\begin{itemize}
				\item $t$: A given time step
				\item $x = \begin{bmatrix} r \\ u \\ v \\ \lambda_r \\ \lambda_u \\ \lambda_v \end{bmatrix}$: the state of the dynamic system vector
				
			\end{itemize}
			\Function{$dynpol$}{$t$, $x$}
			\State $\mu = \text{param}.\mu_{\text{body}}$
			\State $T = \text{param}.T$;
			\State $g0_{\text{Isp}} = \text{param}.g0_{\text{Isp}}$
			\State $m0 = \text{param}.m0$;
			\State $t0 = \text{param}.t0$;
			\State $dm = -\frac{T}{g0_{\text{Isp}}}$
			\State $\phi = \arctan\left(\frac{x(5)}{x(6)}\right)$
			\State $dx = x$;
			\State $dx(1) = x(2)$;
			\State $dx(2) = \frac{x(3)^2}{x(1)} - \frac{\mu}{x(1)^2} + \frac{T}{m0 + dm \times t} \times \sin(\phi)$
			\State $dx(3) = -\frac{x(2) \times x(3)}{x(1)} + \frac{T}{m0 + dm \times t} \times \cos(\phi)$
			\State $dx(4) = -x(5) \times \left(-\frac{x(3)^2}{x(1)^2} + 2 \times \frac{\mu}{x(1)^3}\right) - x(6) \times \frac{x(2) \times x(3)}{x(1)^2}$
			\State $dx(5) = -x(4) + x(6) \times \frac{x(3)}{x(1)};$
			\State $dx(6) = -x(5) \times (2 \times \frac{x(3)}{x(1)}) + x(6) \times \frac{x(2)}{x(1)};$
			\EndFunction
		\end{algorithmic}
	\end{algorithm}
	
	\begin{algorithm}[H]
		\caption{Function \texttt{gnultmin}}
		\begin{algorithmic}
			\State \textbf{Input:}
			\begin{itemize}
				\item $p = \begin{bmatrix} t_f \\ \lambda_{r_0} \\ \lambda_{u_0} \\ \lambda_{v_0} \end{bmatrix}$: the state of the dynamic system vector
			\end{itemize}
			\Function{gnmt}{$p$}
			\State $t0 = \text{param}.t0;$
			\State $r0 = \text{param}.x_{t0}(1);$
			\State $u0 = \text{param}.x_{t0}(2);$
			\State $v0 = \text{param}.x_{t0}(3);$
			\State $rf = \text{param}.x_{tf}(1);$
			\State $uf = \text{param}.x_{tf}(2);$
			\State $vf = \text{param}.x_{tf}(3);$
			\State $y0 = [r0 \quad u0 \quad v0 \quad p(2) \quad p(3) \quad p(4)];$
			\State $tf = p(1);$
			\State $\text{atol} = 1e-10;$
			\State $\text{rtol} = 1e-10;$
			\State $y = \text{ode}("rk", y0, t0, tf, \text{dynpol}, \text{atol}, \text{rtol});$
			\State $gnmt = p;$
			\State $gnmt(1) = y(1) - rf;$
			\State $gnmt(2) = y(2) - uf;$
			\State $gnmt(3) = y(3) - \sqrt{\text{param}.\mu_{\text{body}} / rf};$
			\State $gnmt(4) = p(2)^2 + p(3)^2 + p(4)^2 - 1;$
			\EndFunction
		\end{algorithmic}
	\end{algorithm}
	
	The function \texttt{gnultmin()} takes \( p \) as a parameter and returns the column vector \( gnmt \).
	
	Absolute and relative tolerances of \( 10^{-10} \) have been chosen for solving the differential equation using the function \texttt{ode()}. If the order of magnitude of the relative tolerance is smaller than \( 10^{-10} \), Scilab displays a warning message and increases this tolerance to have an order of magnitude of \( 10^{-10} \).
	
	We aim to achieve better accuracy than that of the solver.
	
	\begin{algorithm}[H]
		\caption{Control History Calculation}
		\begin{algorithmic}[1]
			\For{$i = 1$ \textbf{to} $6$}
			\State $param.T \gets \frac{T(i)}{param.FU}$
			\State $[x, v, info] \gets \text{fsolve}(p0, gnultmin)$
			\State $\text{disp}(info)$
			\State $\text{disp}(v)$
			\State $InitialCond \gets [\text{param.x\_t0(1)}; \text{param.x\_t0(2)}; \text{param.x\_t0(3)}; x(2); x(3); x(4)]$
			\State $t \gets 0:0.01:x(1)$
			\State $sol \gets \text{ode}(InitialCond, param.t0, t, dynpol)$
			\State $timeInDays \gets t \times \frac{param.TU}{86400}$
			\State $AngleInDegrees \gets \text{atan}(sol(5,:), sol(6,:)) \times \frac{360}{2\pi}$
			\State $\text{figure}$
			\State $\text{plot}(timeInDays, AngleInDegrees)$
			\State $\text{xlabel}('Time (days)')$
			\State $\text{ylabel}('Control in deg')$
			\State $powerInNewton \gets param.T \times param.FU$
			\State $\text{xgrid}$
			\State $h \gets \text{gca}()$
			\State $h.\text{background} \gets \text{color}('white')$
			\State $\text{title}('Control history for ' + \text{string}(powerInNewton) + 'N')$
			\EndFor
		\end{algorithmic}
	\end{algorithm}
	
	
	Each thrust (in Newtons) is contained in \( T \). We proceed with increasing thrusts, starting with the lowest one (here 0.1), so the thrust duration decreases. The very first time the function \( gnultmin() \) is passed as a parameter to the function \( fsolve() \), i.e., to calculate the trajectory with the first thrust, we use \( p_0 \) for the initialization of \( gnultmin() \). The result of the first use of \( fsolve() \) called \( p_1 \) will serve for the second time when we pass \( gnultmin() \) as a parameter to \( fsolve() \), hence for calculating the trajectory with the second thrust. Thus, for each thrust except the first one, we use the result of \( fsolve() \) from the previous thrust. With \( p_0 \), we can then obtain \( p_i \) for all \( i \in [1; 6] \).
	
	\begin{figure}[htbp]
		\centering
		\includegraphics[width=\textwidth]{figure0.png}
		\caption{thrusts of 0.1N}
		\label{fig:0.1}
	\end{figure}
	
	\begin{figure}[htbp]
		\centering
		\includegraphics[width=\textwidth]{figure1.png}
		\caption{thrusts of 0.2N}
		\label{fig:0.2}
	\end{figure}
	
	\begin{figure}[htbp]
		\centering
		\includegraphics[width=0.7\textwidth]{figure2.png}
		\caption{thrusts of 0.3N}
		\label{fig:0.3}
	\end{figure}
	
	\begin{figure}[htbp]
		\centering
		\includegraphics[width=0.7\textwidth]{figure3.png}
		\caption{thrusts of 0.4N}
		\label{fig:0.4}
	\end{figure}
	
	\begin{figure}[htbp]
		\centering
		\includegraphics[width=0.7\textwidth]{figure4.png}
		\caption{thrusts of 0.5N}
		\label{fig:0.5}
	\end{figure}
	
	\begin{figure}[htbp]
		\centering
		\includegraphics[width=0.7\textwidth]{figure5.png}
		\caption{thrusts of 0.6N}
		\label{fig:0.6}
	\end{figure}
	
	
\end{document}